% !TeX program = lualatex
%======================================================================
% DOCUMENT SETUP -- compile with recent LuaLaTeX
%======================================================================
\DocumentMetadata{
  tagging=on,
  tagging-setup={math/setup=mathml-SE},
  pdfstandard=ua-2,
  lang=en-US
}
\documentclass[11pt,letterpaper]{article}
\usepackage[widemargin]{notes}
% \usepackage{notes}

%----------------------------------------------------------------------
% PDF PROPERTIES
%----------------------------------------------------------------------
\hypersetup{
  pdflang={en-US},
  pdftitle={\class{} - Day 41: A Historical Review},
  pdfauthor={\instructor},
  pdfsubject={\coursename}
}

\begin{document}

\lesson[George F. Simmons, Differential Equations with
Applications and Historical Notes, 3rd edition]
  {Day 41}{A Historical Review}

%======================================================================
% OPENING
%======================================================================

The goal of this lecture is to connect the major ideas
from our five chapters and place them in historical
context. We follow the order of the course, not a strict
timeline: physical questions led to equations, the search
for formulas led to new methods, and difficult equations
motivated studying solutions without explicit formulas.

%%%%%%%%%%%%%%%%%%%%%%%%%%%%%%%%%%%%%%%%%%%%%%%%%%%%%%%%%%%%%%%%%%%%%%%%
\section{Chapter 1: First-Order Equations}

\subsection{From Calculus to Differential Equations}

\textbf{Main idea: A rate law turns a physical assumption into a differential equation.}

In the late seventeenth century, Isaac Newton and
Gottfried Wilhelm Leibniz independently developed
calculus. Derivatives gave mathematicians a language
for instantaneous change. Differential equations
asked a related question: if we know how a quantity
changes, can we determine the quantity itself?

Some of the earliest problems concerned moving bodies
and curves defined by their tangents. In 1701, Newton
published work on the cooling of hot objects; a modern
idealization is $T'=-k(T-T_a)$. Later, Thomas Malthus
argued in 1798 that unchecked population growth could
be geometric. The equation below is a modern formulation
of proportional growth, not an equation Malthus wrote.

% Historical sources: U. Besson, European Journal of Physics (2011),
% DOI 10.1088/0143-0807/32/2/007 (Newton's 1701 cooling work).
% MacTutor biography of Verhulst and historical work on Malthus (1798).
For example, suppose a population grows at a rate
proportional to its size:
\[
P'=rP.
\]
Separation of variables gives the family
\[
P(t)=Ce^{rt}.
\]
The initial condition $P(0)=P_0$ selects one member:
\[
\boxed{P(t)=P_0e^{rt}.}
\]

The differential equation describes the growth rule.
The initial condition tells us where that growth begins.

\textbf{Segue: A rate law gives us an equation. How can we recover the unknown function?}

%----------------------------------------------------------------------
\subsection{Finding a Method That Fits}

\textbf{Main idea: Recognize a useful structure before choosing an exact solution method.}

Leibniz and the Bernoulli brothers, Jacob and Johann,
were central figures in the early development of
differential equations. Many early problems involved
curves, motion, and geometry. Solving them required
recognizing useful forms and making clever substitutions.
Leibniz treated first-order linear equations in 1694;
Johann Bernoulli and, later, Euler helped turn such
individual insights into more systematic methods.
Our integrating-factor formula condenses that history
into one reliable procedure.

% Historical source: Mathematical Association of America,
% "Solving First-Order Linear Differential Equations:
% Three Mini-Primary Source Projects" (Leibniz, Bernoulli, Euler).

That is also how we approached first-order equations.
Before calculating, we identified the structure.

A \textbf{separable equation} has the form
\[
x'=f(t)g(x).
\]
We separate the variables and integrate. When dividing
by $g(x)$, we must also check whether $g(x)=0$ gives
equilibrium solutions.

A \textbf{first-order linear equation} has the form
\[
x'+p(t)x=q(t).
\]
The integrating factor
\[
\mu(t)=e^{\int p(t)\,dt}
\]
turns the equation into
\[
(\mu x)'=\mu q.
\]

Both methods make the equation easier to integrate,
but they use different features of its structure.

\begin{example}{
Classify both equations, then solve the linear one.
\[
x'=t(1+x^2),
\qquad
x'+2x=t.
\]
}
The first is separable. For the second, $\mu=e^{2t}$ gives
\[
(e^{2t}x)'=te^{2t},\qquad
e^{2t}x=e^{2t}\left(\frac t2-\frac14\right)+C.
\]
Hence
\[
\boxed{x(t)=\frac t2-\frac14+Ce^{-2t}.}
\]
\end{example}


%----------------------------------------------------------------------

\subsection{Existence, Uniqueness, and Behavior}

\textbf{Main idea: Existence, uniqueness, and long-term behavior are separate from finding a formula.}

Finding a formula is only one part of understanding
a differential equation. We also asked:
\begin{itemize}
    \item Does a solution exist?
    \item Does the initial condition determine a
    unique solution?
    \item How does the solution behave over time?
\end{itemize}

During the nineteenth century, Cauchy advanced
existence arguments, while Lipschitz, Peano, and
Picard clarified the assumptions behind existence
and uniqueness. In particular, continuity alone can
give local existence without guaranteeing uniqueness.
This separated the question of whether a solution exists
from whether we can write it down.

% History: Cauchy's work in the 1820s; Peano's existence
% results in the late 1880s; Picard's iterations around 1890.
% See Trefethen et al., Exploring ODEs, historical notes,
% https://people.maths.ox.ac.uk/trefethen/Exploring.pdf

Recall a sufficient condition from our course:
if $f$ and $f_x$ are continuous near $(t_0,x_0)$,
then the IVP
\[
x'=f(t,x),\qquad x(t_0)=x_0
\]
has a unique solution on some interval around $t_0$.
This is a local conclusion; it does not automatically
guarantee a solution for all future time.

We also studied autonomous equations using phase lines.
For example, Pierre-Fran\c{c}ois Verhulst proposed a
logistic population model in 1838 to account for
limits on unchecked growth:
\[
P'=rP\left(1-\frac{P}{K}\right),
\qquad r,K>0.
\]
The equilibria are $P=0$ and $P=K$. Without solving:
\[
\begin{aligned}
0<P<K &\quad\Longrightarrow\quad P'>0,\\
P>K   &\quad\Longrightarrow\quad P'<0.
\end{aligned}
\]
Positive populations move toward the carrying
capacity $K$.

We could find an explicit solution, but the signs
already tell us the main long-term behavior.

\textbf{Segue:} A growth law involves the first
derivative. Newton's force law involves acceleration,
so mechanics naturally leads to second-order equations.

% BLUE TRANSITION BOX: From Rates to Forces
% Growth naturally involves a first derivative.
% Acceleration naturally involves a second derivative.

%%%%%%%%%%%%%%%%%%%%%%%%%%%%%%%%%%%%%%%%%%%%%%%%%%%%%%%%%%%%%%%%%%%%%%%%
\section{Chapter 2: Second-Order Equations}

% SIDE NOTE: 9:12 | Forces Produce Motion
% a) Newton's law connects force and acceleration.
% b) Identify mass, damping, stiffness, and forcing.
% c) The same equation can describe different physical systems.

\subsection{Newton's Law and the Spring Model}

\textbf{Main idea: Force laws naturally give equations involving acceleration.}


Newton's laws of motion, published in the
\textit{Principia} in 1687, connected forces to
changes in motion. Together with Hooke's earlier
work on elastic restoring forces, they provide
building blocks for the modern spring model.
For constant mass,
\[
F=mx''.
\]

For a spring-mass system, the forces lead to
\[
\boxed{mx''+\gamma x'+kx=F(t).}
\]
Here,
\begin{itemize}
    \item $m$ is the mass;
    \item $\gamma$ is the damping coefficient;
    \item $k$ is the spring constant;
    \item $F(t)$ is the applied external force.
\end{itemize}

The restoring force is $-kx$, and the damping force
is $-\gamma x'$. Moving these terms to the left
gives the equation above.

The same mathematical structure appears in electrical
circuits. Studying one equation can therefore help us
understand several kinds of physical systems.

\textbf{Segue:} Given a force law, can we predict the motion by finding an explicit formula?

%----------------------------------------------------------------------
\subsection{Euler and Exponential Solutions}
\textbf{Main: Exponential trial solutions transform linear differential equations into algebraic problems. }

Leonhard Euler helped organize and extend the methods
of differential equations during the eighteenth century.
In the 1730s and 1740s, Euler and Johann Bernoulli
studied higher-order linear equations; Euler's published
treatment connected constant-coefficient differential
equations to algebraic ones.

% Historical source: MAA, "Solving Linear Higher Order
% Differential Equations with Euler and Johann Bernoulli,"
% discussing their correspondence (1739) and Euler's 1743 work.

The key observation is familiar:
differentiating an exponential preserves its form.
For
\[
mx''+\gamma x'+kx=0,
\]
try
\[
x=e^{\lambda t}.
\]
Substitution gives
\[
(m\lambda^2+\gamma\lambda+k)e^{\lambda t}=0.
\]
Since the exponential is never zero,
\[
\boxed{m\lambda^2+\gamma\lambda+k=0.}
\]

This is the \textbf{characteristic equation}.
Its roots determine the forms of the solutions:
\begin{itemize}
    \item Distinct real roots give two exponential terms.
    \item A repeated real root introduces a factor of $t$.
    \item Complex roots give sine and cosine terms
    multiplied by an exponential.
\end{itemize}

The connection between complex exponentials and
oscillations is expressed by Euler's formula:
\[
e^{i\beta t}=\cos\beta t+i\sin\beta t.
\]

\begin{example}{
Compare the general solutions of
\[
x''+4x=0,
\qquad
x''+2x'+4x=0.
\]
}
The characteristic roots are $\lambda=\pm2i$
and $\lambda=-1\pm i\sqrt3$, respectively. Thus
\[
x_1=c_1\cos2t+c_2\sin2t,
\]
\[
x_2=e^{-t}\bigl(c_1\cos(\sqrt3 t)
                   +c_2\sin(\sqrt3 t)\bigr).
\]
The first oscillates with constant amplitude;
the second has a decaying amplitude.
\end{example}

%----------------------------------------------------------------------
% SIDE NOTE: 9:18 | Natural Motion and Forcing
% a) Write x = x_h + x_p.
% b) Recall why a duplicate trial fails.
% c) Connect the repaired trial to resonance.

\subsection{Superposition, Forcing, and Resonance}
\textbf{Main idea: Superposition separates natural motion from the response to a force.}

The characteristic equation describes a spring
left to move on its own, with $F(t)=0$. But a
periodic push changes the model to
$mx''+\gamma x'+kx=F(t)$. What must we add to
the natural motion to account for that push?
Linearity supplies the answer.

For a nonhomogeneous equation, we used
\[
\boxed{x=x_h+x_p.}
\]
The homogeneous solution contributes zero to the
forcing, while the particular solution contributes
$F(t)$. Their sum therefore solves the required equation.

The method of undetermined coefficients finds $x_p$
by choosing a trial and determining its coefficients.
Our failed trials taught us something important:
a trial already belonging to the homogeneous family
cannot produce a nonzero forcing term.

For example,
\[
x''+4x=\cos2t
\]
has homogeneous solution
\[
x_h=c_1\cos2t+c_2\sin2t.
\]
The usual sine/cosine trial is a duplicate.
Multiplying the trial group by $t$ leads to
\[
x_p=\frac{t}{4}\sin2t.
\]

The growing factor $t$ also has a physical meaning.
In this ideal undamped model, forcing at the natural
frequency produces \textbf{resonance}: oscillations
with a growing amplitude.

For damped systems, we also distinguished the
decaying \textbf{transient response} from the
long-term response driven by the external force.

\textbf{Segue:} Polynomial, exponential, and sinusoidal
forcing work well with MOUC. What if a voltage
switches on suddenly or an input acts almost
instantaneously?

%%%%%%%%%%%%%%%%%%%%%%%%%%%%%%%%%%%%%%%%%%%%%%%%%%%%%%%%%%%%%%%%%%%%%%%%
\section{Chapter 3: Laplace Transforms}



\subsection{Transform, Solve, and Invert}
\textbf{Main idea: Transforms turn differentiation into algebra while retaining the initial data.}

Pierre-Simon Laplace used related integral methods
in probability and mathematical physics in the late
eighteenth and early nineteenth centuries.
The streamlined transform technique we use today
was systematized much later; Laplace did not
work from our modern table of transforms.

% History: C. C. Gillispie, Pierre-Simon Laplace, 1749--1827
% (Princeton University Press); see discussion of the
% subsequent development of the Laplace transform.

The \textbf{Laplace transform} assigns a new function
to $x(t)$:
\[
X(s)=\mathcal{L}\{x(t)\}
=\int_0^\infty e^{-st}x(t)\,dt,
\]
for values of $s$ where the integral converges.

Our most useful identities were
\[
\mathcal{L}\{x'\}=sX(s)-x(0)
\]
and
\[
\mathcal{L}\{x''\}
=s^2X(s)-sx(0)-x'(0).
\]
Derivatives become algebraic expressions involving
$X(s)$, and the initial conditions enter directly.

\begin{example}{
Solve this IVP using the Laplace transform.
\[
x'+2x=1,\qquad x(0)=0.
\]
}
The equation
\[
(s+2)X=\frac1s.
\]
Therefore,
\[
X=\frac{1}{s(s+2)}
=\frac12\left(\frac1s-\frac1{s+2}\right),
\]
so
\[
\boxed{x(t)=\frac12(1-e^{-2t}).}
\]
An integrating factor would also solve this problem.
The value of the transform method becomes especially
clear when the input is piecewise or impulsive.
\end{example}

%----------------------------------------------------------------------
\subsection{Switches and Impulses}

\textbf{Main idea: Steps and impulses represent real switching and short-duration inputs.}

As telegraph and telephone networks expanded in the
late nineteenth century, Oliver Heaviside studied
electrical signals propagating through wires.
His operational calculus replaced differentiation
by algebraic manipulation, making transient circuit
problems easier to tackle, though the rules still
needed rigorous justification.

% History: MacTutor biography of Oliver Heaviside;
% "On Heaviside's contributions to transmission line theory"
% (Philosophical Transactions A, 2018).

Heaviside's original operational calculus was not
identical to the Laplace-transform method we learned.
However, both illustrate the same useful ambition:
turn a differential equation into an algebraic problem.

The \textbf{unit step function}
\[
H(t-a)
\]
describes an input that switches on at $t=a$.

The \textbf{Dirac delta}
\[
\delta(t-a)
\]
represents an idealized instantaneous impulse.
It is not an ordinary function with a finite spike.
The notation became widespread through Dirac's
work in quantum mechanics in the twentieth century;
impulse-like mathematical ideas are older than Dirac.

% History: Stanford Encyclopedia of Philosophy,
% "Quantum Theory and Mathematical Rigor," section on
% Dirac's delta and his 1930 textbook.

Their Laplace transforms are
\[
\mathcal{L}\{H(t-a)\}=\frac{e^{-as}}{s},
\qquad
\mathcal{L}\{\delta(t-a)\}=e^{-as}.
\]

For example, in a series RLC circuit the charge $q$
and current $i=q'$ satisfy
\[
Lq''+Rq'+\frac1Cq=V(t).
\]
A voltage $V_0H(t-a)$ models switching a source on
at $t=a$. An ideal voltage impulse $J\delta(t-a)$
models an extremely short pulse; it changes the
inductor current by $J/L$ in the ideal circuit.

% OPTIONAL worked example if short on lecture time.
\begin{example}{
A series RL circuit has $L=1$, $R=2$, and $i(0)=0$.
A 6-volt source switches on at $t=5$. Find $i(t)$.
}
The equation is $i'+2i=6H(t-5)$.
Transforming gives
\[
(s+2)I(s)=\frac{6e^{-5s}}s,\qquad
I(s)=3e^{-5s}\left(\frac1s-\frac1{s+2}\right).
\]
Using the second shifting theorem,
\[
\boxed{i(t)=3\bigl(1-e^{-2(t-5)}\bigr)H(t-5).}
\]
The current is zero until the switch closes.
\end{example}

%%%%%%%%%%%%%%%%%%%%%%%%%%%%%%%%%%%%%%%%%%%%%%%%%%%%%%%%%%%%%%%%%%%%%%%%
\section{Chapter 4: Linear Systems}

% SIDE NOTE: 9:30 | Find the Basic Motions
% a) Interacting objects lead to coupled equations.
% b) A second-order equation is also a first-order system.
% c) The exponential trial produces an eigenvalue problem.

\subsection{From One Quantity to Several}

\textbf{Main idea: Interacting quantities must be modeled together as a system.}

Problems involving several interacting objects
helped drive the development of differential equations.
The eighteenth-century study of vibrations involved
mathematicians such as Euler, Daniel Bernoulli,
d'Alembert, and Lagrange. Their work on small
oscillations and normal modes asked how coupled bodies
move together. Celestial mechanics likewise required
studying bodies that exert forces on one another.

The matrix notation below is modern shorthand for
these coupled equations, not how eighteenth-century
mathematicians originally wrote them:
\[
\mathbf{x}'=A\mathbf{x}.
\]

Even a single second-order equation can be rewritten
as a first-order system. For example,
\[
x''+4x=0
\]
becomes, with $y=x'$,
\[
\begin{pmatrix}x\\y\end{pmatrix}'
=
\begin{pmatrix}
0&1\\
-4&0
\end{pmatrix}
\begin{pmatrix}x\\y\end{pmatrix}.
\]
The coordinates now record position and velocity
at the same time.

%----------------------------------------------------------------------
\subsection{The Exponential Idea Returns}

\textbf{Main idea: Exponential modes reduce a linear system to an eigenvalue problem.}

For a scalar equation, we tried $e^{\lambda t}$.
For a system, we tried
\[
\mathbf{x}=e^{\lambda t}\mathbf{v}.
\]
Substituting gives
\[
\lambda e^{\lambda t}\mathbf{v}
=e^{\lambda t}A\mathbf{v},
\]
and therefore
\[
\boxed{A\mathbf{v}=\lambda\mathbf{v}.}
\]

The \textbf{eigenvalue problem} arises naturally
from the search for exponential solutions. It also
formalizes the older physical idea of \emph{normal modes}:
special motions in which the components share a common
time dependence.

For distinct real eigenvalues,
\[
\mathbf{x}(t)
=c_1e^{\lambda_1t}\mathbf{v}_1
+c_2e^{\lambda_2t}\mathbf{v}_2.
\]
Real eigenvectors determine the directions of
linear orbits. The signs of the eigenvalues
determine whether those orbits point toward
or away from the origin.

Complex eigenvalues introduce rotation, while
their real parts determine growth or decay.

%----------------------------------------------------------------------
\subsection{From Formulas to Phase Portraits}

\textbf{Main idea: A phase portrait summarizes whole families of motions without solving each IVP.}

A \textbf{phase portrait} displays representative
trajectories, with arrows indicating increasing time.
It lets us compare many initial conditions in one picture.

We classified the familiar cases:
nodes, saddles, lines of equilibria, centers,
spirals, stars, and degenerate nodes.

\begin{example}{    
Classify the origin for
\[
A_1=
\begin{pmatrix}2&0\\0&-1\end{pmatrix},
\qquad
A_2=
\begin{pmatrix}-1&0\\0&-2\end{pmatrix},
\qquad
A_3=
\begin{pmatrix}0&1\\-4&0\end{pmatrix}.
\]
}
The eigenvalues are $2,-1$ for $A_1$,
$-1,-2$ for $A_2$, and $\pm2i$ for $A_3$.
Thus these give a saddle, a stable node, and a
center, respectively.

Notice the connection to Chapter 2. The matrix
$A_3$ describes the same undamped oscillator
$x''+4x=0$. Its eigenvalues are the characteristic
roots we found earlier.
\end{example}

% BLUE TRANSITION BOX: Understand Without a Formula
% Linear systems let us connect formulas and pictures.
% For nonlinear systems, the picture may be our main source of information.

%%%%%%%%%%%%%%%%%%%%%%%%%%%%%%%%%%%%%%%%%%%%%%%%%%%%%%%%%%%%%%%%%%%%%%%%
\section{Chapter 5: Nonlinear Systems}

% SIDE NOTE: 9:38 | A Different Kind of Answer
% a) Poincare studied the behavior of solutions.
% b) Find equilibria and linearize nearby.
% c) Local information does not determine the full portrait.

\subsection{Poincar\'e and Qualitative Questions}

\textbf{Main idea: Geometry and stability can reveal behavior when explicit solutions are unavailable.}

Poincar\'e began a systematic study of trajectories
defined by differential equations in the early 1880s.
His later work on celestial mechanics, especially
the three-body problem, showed how powerful that
geometric viewpoint could be. Newton's laws gave
the equations of interacting bodies, but not a
general explicit formula for their motions.
Poincar\'e instead studied the geometry of possible
trajectories, periodic motion, and stability.

% History: Poincare's "curves defined by differential
% equations" memoirs (1880s) preceded his three-body work;
% see the history of the three-body problem (1889 prize
% and Methodes nouvelles, 1892--1899).

Instead of asking only for a solution formula,
one could ask:
\begin{itemize}
    \item Does the motion remain bounded?
    \item Does it approach an equilibrium?
    \item Does it repeat?
    \item What happens after a small change
    in the initial condition?
\end{itemize}

This viewpoint brings us back to the phase lines
from Chapter 1, but now several quantities
can change together.

%----------------------------------------------------------------------
\subsection{Linearization Connects Two Chapters}

\textbf{Main idea: Nearby nonlinear motion can often be approximated by a familiar linear system.}

For
\[
x'=f(x,y),\qquad y'=g(x,y),
\]
an \textbf{equilibrium} $(x^*,y^*)$ satisfies
\[
f(x^*,y^*)=0,\qquad g(x^*,y^*)=0.
\]

Near that point, we use the \textbf{Jacobian matrix}
\[
J(x^*,y^*)=
\begin{pmatrix}
f_x(x^*,y^*)&f_y(x^*,y^*)\\
g_x(x^*,y^*)&g_y(x^*,y^*)
\end{pmatrix}.
\]
With
\[
u=x-x^*,\qquad v=y-y^*,
\]
the local linear approximation is
\[
\begin{pmatrix}u\\v\end{pmatrix}'
=
J(x^*,y^*)
\begin{pmatrix}u\\v\end{pmatrix}.
\]

This is why our work on linear systems was so useful:
it gives us a way to study a nonlinear system
near each equilibrium.

However, linearization is local. It does not
automatically tell us how trajectories connect
far from the equilibrium.

%----------------------------------------------------------------------
% SIDE NOTE: 9:43 | From Mechanics to Populations
% a) Verhulst modeled one population.
% b) Lotka and Volterra studied interacting populations.
% c) Read the terms before doing any calculations.

\subsection{Population Dynamics}

\textbf{Main idea: Model assumptions about interactions can predict population cycles.}

Differential equations also became tools for studying
living systems. Verhulst's logistic model described
a single population with limited resources.
In the 1920s, Alfred Lotka studied interacting
populations, inspired partly by physical chemistry.
Independently, Vito Volterra investigated predator--prey
cycles after a question about Adriatic fishing data
raised by the biologist Umberto d'Ancona.

% History: S. Kingsland, "Alfred J. Lotka and the origins
% of theoretical population ecology," PNAS (2015);
% A Short History of Mathematical Population Dynamics,
% chapter "Lotka, Volterra and the predator--prey system".

Our predator--prey model was
\[
\begin{dcases}
x'=rx-\alpha xy,\\
y'=-ky+\beta xy,
\end{dcases}
\qquad r,k,\alpha,\beta>0.
\]

Each term expresses an assumption:
\begin{itemize}
    \item $rx$: prey grow in the absence of predators.
    \item $-ky$: predators decline in the absence of prey.
    \item $-\alpha xy$: encounters reduce the prey population.
    \item $\beta xy$: food supports predator growth.
\end{itemize}

The two interaction coefficients need not be equal.
Prey consumed and predator growth are different
quantities.

We interpreted the terms, found equilibria, and
studied trajectories without needing explicit
formulas for both populations. This is a useful
example of Poincar\'e's broader viewpoint:
understanding a solution does not always require
writing it down.

% OPTIONAL: algebraic recap of nonlinear linearization.
\begin{example}{
Find the positive coexistence equilibrium of the
predator--prey system and classify its linearization.
}
Setting $x'=y'=0$ gives
$(x^*,y^*)=(k/\beta,r/\alpha)$. At this point,
\[
J(x^*,y^*)=
\begin{pmatrix}
0&-\alpha k/\beta\\
\beta r/\alpha&0
\end{pmatrix},
\qquad \lambda^2=-rk.
\]
The linearized system has a center. Purely imaginary
eigenvalues alone do not prove that a \emph{nonlinear}
equilibrium is a center; this model has an additional
conserved quantity that establishes closed orbits.
\end{example}

\end{document}